
La realización de este proyecto ha seguido una metodología colaborativa para el diseño conceptual y lógico, dividiendo posteriormente la implementación técnica y el desarrollo en Oracle APEX por módulos funcionales. A continuación, se detalla la intervención de cada miembro del equipo:

\begin{enumerate}
    \item \textbf{Diseño Conceptual y Lógico (Intervención Conjunta):} \\
    Todos los miembros del grupo \textit{(Ismael Sallami Moreno, Jesús Rodríguez González, Fernando José Gracia Choin, Sergio Calvo González y Javier Niño Sánchez)} participaron activamente en la fase de análisis inicial. Esta etapa incluyó:
    \begin{itemize}
        \item La toma de decisiones respecto a la fusión de tablas para optimizar el modelo, destacando la fusión de las tablas 3, 6, 14 y 15 en la tabla \texttt{PUBLICACION} y la fusión de las tablas 5 y 12 en la estructura de \texttt{MENSAJE}.
    \end{itemize}

    \item \textbf{Implementación del Sistema para el Login:} \\
    \textbf{Ismael Sallami Moreno} fueron los encargados del desarrollo de la tabla \texttt{USUARIO} (tabla 1). Su trabajo abarcó la creación de la Página 2 (\textit{Modificar Usuario}), la configuración de los esquemas de autenticación personalizados (\textit{Login Tabla Usuario}) y el diseño de la Página 9999 (\textit{Login}).

    \item \textbf{Gestión de Contenido y Publicaciones:} \\
    \textbf{Jesús Rodríguez González} asumió la responsabilidad de la tabla \texttt{PUBLICACION} (tabla 2). Implementó la lógica para la gestión de contenido multimedia y desarrolló la Página 3 (\textit{Crear publicaciones}).

    \item \textbf{Módulo de Publicidad y Fusiones Complejas:} \\
    \textbf{Fernando José Gracia Choin} se ocupó de la implementación de la tabla resultante \texttt{PUBLICACION}, producto de la fusión estratégica de las tablas 3, 6, 14 y 15. Desarrolló la Página 11 (\textit{Asignar Anuncio a Publicación}).

    \item \textbf{Sistema de Mensajería y Comunicaciones:} \\
    \textbf{Sergio Calvo González} desarrolló el módulo de comunicación asíncrona, implementando las tablas \texttt{MENSAJE} y \texttt{CONVERSA} (fusión de tablas 5 y 12). Además, finalizó con la creación de la Página 12 (\textit{Enviar Mensajes}) y la gestión de conversaciones.

    \item \textbf{Analítica, Reportes y Redacción de la Memoria:} \\
    \textbf{Javier Niño Sánchez} implementó la Página 4 (\textit{Listar Tendencias}) y la lógica PL/SQL. Asimismo, ha sido el encargado principal de la redacción, estructuración y compilación final de la memoria del proyecto, integrando las documentaciones técnicas aportadas por cada miembro.
\end{enumerate}

\section*{Cambios Realizados}

Durante el desarrollo, se han identificado y realizado los siguientes cambios relevantes respecto a la práctica anterior:

\begin{itemize}
    \item El proyecto parte de la base de la Práctica 2, sobre la cual ya se habían realizado modificaciones previas antes de iniciar este desarrollo.
    \item En la funcionalidad de \textit{Enviar Mensajes}, no se ha tenido en cuenta la restricción de que los usuarios deban ser amigos para poder intercambiar mensajes. Esto supone una diferencia respecto a posibles requisitos de control de acceso más estrictos.
\end{itemize}